\capitulo{7}{Conclusiones y líneas de trabajo futuras}

Todo proyecto debe incluir las conclusiones que se derivan de su desarrollo, destacando los logros alcanzados y los aprendizajes obtenidos. Este capítulo final presenta una síntesis de los resultados del proyecto, tanto funcionales como técnicos. Además, se ofrece un informe crítico detallado que expone las limitaciones del prototipo actual y delinea una serie de líneas de trabajo futuras, que permiten trazar una hoja de ruta para la evolución y expansión del laboratorio virtual de programación y robótica.

\section{Conclusiones}

El presente Trabajo Fin de Grado ha culminado con el desarrollo de un prototipo funcional de un laboratorio virtual para la enseñanza de la programación y la robótica, implementado en el motor Unity 6 con C\#. El proyecto ha logrado con éxito el objetivo principal de crear un entorno de aprendizaje híbrido, que combina una interfaz visualmente estructurada con un sistema de programación textual basado en un lenguaje de dominio específico (DSL) propio. Esta solución ha demostrado la viabilidad técnica de un modelo pedagógico de transición, donde los usuarios aprenden los fundamentos de la algoritmia y el control de un robot 3D sin la curva de aprendizaje inicial de la sintaxis compleja, pero practicando la escritura de código.


Los principales logros y aprendizajes del proyecto incluyen:

\begin{itemize}
	\item \textbf{Implementación de un motor de programación visual viable:} Se ha replicado con éxito la lógica de interacción visual de Scratch (arrastre, conexión, disposición de categorías) y se ha establecido un robusto motor de ejecución directa de bloques, basado en la adaptación de UBlockly. Este motor ha demostrado la capacidad de interpretar y ejecutar una secuencia lineal de bloques (\texttt{motion\_movesteps}, \texttt{event\_whenflagclicked}) y de gestionar corrutinas para la animación fluida del robot.
	\item \textbf{Diseño de una arquitectura MVC efectiva:} La adopción del patrón Modelo-Vista-Controlador ha permitido una clara separación de intereses entre la lógica del programa (modelos de UBlockly como \texttt{Block} y \texttt{Connection}), la representación visual (\texttt{BlockView} en Unity UI), y el control de interacciones (\texttt{BlockController}}). Esto garantiza la mantenibilidad y futura escalabilidad del sistema.
	\item \textbf{Integración exitosa de componentes:} Se ha logrado una integración fluida de los elementos de UBlockly (modelo de bloques, sistema de ejecución, serialización XML) con el entorno de Unity (UI Canvas, objetos 3D, corrutinas). Este proceso ha implicado resolver desafíos complejos relacionados con la sincronización de estados y posiciones entre los modelos lógicos y las vistas visuales de Unity.
	\item \textbf{Validación de un flujo completo de simulación:} Se ha verificado que un programa básico de bloques (ej. al presionar bandera verde, mover 10 pasos) se traduce correctamente en una acción visible y controlada por el robot virtual en el entorno 3D, demostrando la validez del concepto central del laboratorio.
	\item \textbf{Gestión de proyectos en un entorno agile:} La aplicación de la metodología SCRUM, utilizando Zube y GitHub, ha sido fundamental para gestionar la complejidad, priorizar tareas y adaptarse a los retos técnicos emergentes durante el desarrollo del prototipo.
\end{itemize}

\section{Líneas de trabajo futuras}

El presente proyecto establece una base sólida sobre la que se pueden desarrollar mejoras y ampliaciones significativas para que el laboratorio virtual evolucione hacia una solución educativa integral. A continuación, se presentan algunas de las posibles líneas de trabajo futuras:

\subsection{Extensión y diversificación del sistema de bloques}
Inicialmente, el prototipo implementa un conjunto limitado de bloques para demostrar su viabilidad. Para una solución completa, se plantea:
\begin{itemize}
	\item \textbf{Ampliación de categorías de bloques:} Incorporar todas las categorías estándar de Scratch (ej. Sonido, Apariencia, Sensores).
	\item \textbf{Implementación de bloques avanzados:} Extender la funcionalidad para incluir estructuras de control más complejas como bucles, condicionales, funciones (procedimientos y bloques My Blocks), listas y operaciones con variables y operadores lógicos/matemáticos avanzados. Esto implicaría la adaptación de los intérpretes (\texttt{Cmdtor}s) y la generación de prefabs visuales para cada nuevo tipo de bloque.
	\item \textbf{Expansión del modelo de eventos:} Desarrollar la capacidad para gestionar y responder a un rango más amplio de eventos, más allá de la bandera verde, permitiendo interacciones más complejas con el entorno virtual.
	\item \textbf{Sistema de depuración visual interactiva:} Crear herramientas visuales en la interfaz que permitan a los estudiantes seguir el flujo de ejecución del programa paso a paso, visualizar el estado de las variables y detectar errores en tiempo real, facilitando la comprensión y la resolución de problemas de lógica.
	\item \textbf{Integración de bloques con hardware simulado detallado:} Extender los bloques para representar el comportamiento detallado de motores (velocidad, dirección) y sensores (distancia, color, toque) del LEGO WeDo 2.0.
\end{itemize}

\subsection{Mejora de la simulación 3D e interactividad}
Actualmente, la simulación se centra en el movimiento básico del robot. Para mejorar la experiencia de aprendizaje y la inmersión, se pueden introducir mejoras como:
\begin{itemize}
	\item \textbf{Entornos dinámicos e interactivos:} Diseñar escenarios 3D donde el robot pueda interactuar con obstáculos, objetos móviles o elementos activables (puertas, palancas), y responder a estímulos del entorno simulado (ej. zonas con distinto tipo de terreno que afecten el movimiento, etc.).
	\item \textbf{Simulación de física avanzada:} Integrar el motor de físicas de Unity (\texttt{Rigidbody}, \texttt{Collider}) de forma más compleja para dotar al robot y a los objetos del entorno de un comportamiento físico más realista (ej. colisiones, gravedad, rozamiento), lo que afectaría de manera auténtica la movilidad del robot y sus interacciones.
	\item \textbf{Soporte para múltiples modelos de robots:} Expandir el sistema para incluir y permitir la programación de distintos modelos de robots educativos (además del LEGO WeDo 2.0 básico), como LEGO SPIKE Essential, LEGO SPIKE Prime o LEGO MINDSTORMS EV3, cada uno con sus características y modelos 3D detallados.
	\item \textbf{Realidad Virtual (VR) y Aumentada (AR):} Explorar la portabilidad y adaptación del prototipo a plataformas de Realidad Virtual (VR) o Aumentada (AR), utilizando las capacidades nativas de Unity para ofrecer experiencias aún más inmersivas y atractivas para el aprendizaje.
\end{itemize}

\subsection{Evaluación de la efectividad educativa y usabilidad}
Una limitación inherente a la fase de prototipado es la ausencia de una validación exhaustiva con usuarios reales. Para futuras adaptaciones, es fundamental realizar:
\begin{itemize}
	\item \textbf{Estudios de usabilidad y experiencia de usuario (UX):} Realizar pruebas formales con grupos de estudiantes de diversas edades y niveles de experiencia, empleando metodologías como el test A/B o encuestas de satisfacción, para evaluar la facilidad de uso de la interfaz, la intuición del sistema de bloques y la efectividad general de la plataforma.
	\item \textbf{Análisis del impacto en el aprendizaje:} Diseñar y llevar a cabo estudios longitudinales para medir cómo la plataforma mejora las habilidades de programación y el pensamiento computacional de los estudiantes, en comparación con métodos de enseñanza tradicionales o con otras herramientas existentes.
	\item \textbf{Recopilación de métricas de interacción y comportamiento:} Implementar herramientas de análisis y telemetría para recolectar datos anónimos sobre cómo los estudiantes interactúan con el sistema (tiempo dedicado, bloques utilizados, errores comunes, caminos de solución), lo cual permitiría identificar áreas de mejora y optimizar la experiencia de aprendizaje.
\end{itemize}

\subsection{Integración con hardware físico}
Aunque el prototipo actual se centra en la simulación virtual, el objetivo a largo plazo incluye la interacción con robots reales. Las posibles mejoras en esta línea abarcan:
\begin{itemize}
	\item \textbf{Conexión en tiempo real con LEGO WeDo 2.0 (Física):} Desarrollar un módulo de comunicación (mediante Bluetooth Low Energy) que permita la ejecución directa del código programado con bloques en el robot físico LEGO WeDo 2.0, creando un puente entre la programación virtual y la experiencia tangible.
	\item \textbf{Compatibilidad con otros kits de robótica:} Ampliar el sistema para permitir la programación y el control de otros dispositivos de robótica educativa, como Arduino o micro:bit, a través de interfaces de programación de bajo nivel (APIs o protocolos de comunicación específicos) para esos dispositivos.
\end{itemize}

\subsection{Optimización del rendimiento, portabilidad y localización}
Para maximizar la accesibilidad y el alcance del laboratorio virtual, es esencial abordar aspectos de optimización y portabilidad.
\begin{itemize}
	\item \textbf{Optimización del rendimiento en Unity:} Llevar a cabo una fase de optimización exhaustiva para reducir el consumo de recursos de la aplicación (CPU, GPU, RAM), lo cual permitiría su ejecución más fluida en dispositivos de bajo rendimiento, como tablets básicas o computadoras más antiguas, maximizando la inclusión.
	\item \textbf{Portabilidad a versión web del entorno:} Desarrollar una versión ejecutable del laboratorio virtual directamente en navegadores web (utilizando tecnologías como Unity WebGL), eliminando la necesidad de instalación y facilitando un acceso instantáneo desde cualquier dispositivo con conexión a internet.
	\item \textbf{Localización del Sistema:} Traducir la interfaz de usuario, las descripciones de los bloques y los tutoriales a múltiples idiomas. Esto expandiría el alcance educativo de la plataforma a nivel global, eliminando barreras lingüísticas.
\end{itemize}


\subsection{Reflexiones sobre el proceso de desarrollo y la adaptación}
Una de las lecciones más significativas aprendidas durante este TFG ha sido la gestión de la complejidad técnica al adaptar repositorios de código abierto existentes. El objetivo inicial de replicar completamente el sistema visual de Scratch, tomando UBlockly como base, se reveló como un desafío de una magnitud considerablemente mayor a la anticipada.

El proceso de ingeniería inversa para comprender el funcionamiento interno de UBlockly, una librería con documentación limitada, requirió un gran esfuerzo de investigación y análisis. Lo que inicialmente parecía un problema sutil de implementación de layouts visuales se convirtió en un obstáculo significativo, especialmente en la recreación de bloques de control de flujo complejos (los denominados bloques de tipo C, como \texttt{if/else} o bucles), que requieren una manipulación dinámica de la geometría y la jerarquía de las vistas.

Lejos de ser un fracaso, este desafío representó un punto de inflexión crítico y formativo. Puso de manifiesto la importancia de la evaluación realista de los plazos y del alcance en un proyecto de ingeniería de software. La decisión de pivotar hacia un modelo híbrido no fue un abandono del objetivo, sino una redefinición estratégica para garantizar la entrega de un producto funcional y coherente, centrando el esfuerzo en los aspectos donde se podía aportar mayor valor: el diseño de un lenguaje, su intérprete y la simulación interactiva.

Esta experiencia ha reforzado la conciencia sobre la dificultad de integrar sistemas complejos y ha subrayado la importancia de la modularidad y la arquitectura flexible, principios que guiaron tanto la refactorización inicial a MVC como el diseño final del prototipo. Los conocimientos adquiridos durante este proceso de adaptación y resolución de problemas constituyen uno de los aprendizajes más valiosos del proyecto

\section{Conclusión final}

El desarrollo de este laboratorio virtual representa una trayectoria de ingeniería completa, desde la concepción de una idea ambiciosa hasta la entrega de un prototipo funcional y pragmático. La investigación inicial, centrada en la adaptación de UBlockly, demostró rápidamente la inmensa complejidad de replicar sistemas de programación visual como Scratch, especialmente en la generación dinámica de bloques de control de flujo. Este desafío, lejos de ser un fracaso, se convirtió en una valiosa lección sobre la gestión de la complejidad y la importancia de la adaptación estratégica.

El pivote hacia un enfoque híbrido ha culminado en un prototipo que valida un modelo pedagógico original y técnicamente sólido. La aplicación de los conocimientos del Grado en Ingeniería Informática ha sido fundamental, no solo para implementar el sistema final de intérprete de comandos, sino también para analizar la arquitectura inicial, identificar sus obstáculos y tomar decisiones de diseño informadas. El proyecto resultante es, por tanto, una prueba de las habilidades de ingeniería desarrolladas y un reflejo del compromiso con la creación de herramientas educativas accesibles, en línea con los Objetivos de Desarrollo Sostenible.

Las líneas de trabajo futuras, que incluyen retomar la implementación completa del sistema visual, no se plantean como una tarea pendiente, sino como el siguiente paso lógico en un proyecto de mayor envergadura. Personalmente, este Trabajo Fin de Grado ha sido un aprendizaje profundo en arquitectura de software, desarrollo avanzado en Unity y resolución de problemas complejos. La intención de continuar esta línea de investigación, posiblemente durante un futuro máster en diseño de videojuegos, subraya la pasión y el conocimiento generados por este TFG, que considero un primer paso firme y exitoso.

